\section{\label{I_1_C}Un modèle adapté aux archives audiovisuelles : le manuel de catalogage des images animées de la FIAF } 

L’adaptation du modèle FRBR aux archives audiovisuelles impose une double redéfinition : elle doit prendre en compte d’une part la nature propre du langage cinématographique, et d’autre part, les spécificités des corpus audiovisuels conservés. La modélisation en différentes couches de description prévue par le FRBR est pertinente pour le cinéma, qui partage avec le milieu bibliographique un mode de production industriel, permettant ainsi de distinguer la création intellectuelle de ses copies et diffusions multiples. Toutefois, les niveaux bibliographiques classiques dits WEMI (Work, Expression, Manifestation, Item) restent limitées pour la description métier des spécificités des archives audiovisuelles.  

Premièrement, alors que les bibliothèques conservent principalement des copies identiques d’une même publication, les archives d’images animées sont fréquemment amenées à gérer des objets uniques et rares tels que des \textit{rushes}, des films amateurs, montages inachevés, copies annotées, etc. \footcite{zotero-item-646}. Pour ce type de collections, le niveau Manifestation devient obsolète, en raison du caractère inédit de leur création. Deuxièmement, la distinction du FRBR bibliographique entre l’Œuvre (l’idée de création), et l’Expression (sa réalisation) s’adapte mal au contexte cinématographique. En effet, un film prend véritablement forme au moment du tournage puis du montage, rendant ainsi l’idée d’une existence abstraite de création en amont, caduque. Il parait alors plus pertinent d’appréhender l’œuvre cinématographique moins comme une entité abstraite en amont de son expression que comme une entité concrète ancrée dans un processus d’expression.  

Afin de répondre à ces contraintes imposées par le processus de création cinématographique, la FIAF a élaboré un manuel de catalogage des images animées \footcite{zotero-item-644}. Le manuel adapte le modèle FRBR en le combinant aux règles descriptives RDA (Ressources, Description, Accès) dédiées aux archives, et aux normes de description des œuvres cinématographiques (EN 15744 / EN 15907) \footcite{zotero-item-646}. Le manuel se focalise sur les entités de premier groupe afin d’en repenser l’articulation et de proposer de remplacer le schéma WEMI par le schéma WVMI (Work, Variant, Manifestation, Item).  

En effet la FIAF propose de fusionner le niveau Expression dans celui de l'Œuvre, permettant alors à l’Œuvre de regrouper la description intellectuelle et le processus de création de l’archive cinématographique. Par exemple, si nous prenons le film \textit{L'Aigle à deux têtes} (France, 1948, réalisé par Jean Cocteau), le niveau Oeuvre permettrait d’y décrire le contexte de création ainsi que le contenu. De plus, la FIAF propose une nouvelle entité, celle de Variante, optionnelle, qui se substitue à celle d’Expression. Elle permet de décrire toute version différente d’une œuvre originale dans la mesure où son sens principal n’est pas altéré. Par exemple, une version d’un film qui serait colorisée, censurée, doublée, remontée, etc. Si nous reprenons \textit{L'Aigle à deux têtes}, une variante possible est celle de la restauration du film en 4K en 2023, dont le niveau Variante permettrait de décrire les variations portées sur l’œuvre. La FIAF conserve par ailleurs le niveau de Manifestation, toutefois redéfini afin de pouvoir décrire la matérialisation physique ou numérique d’une Œuvre ou de sa Variante sur un support, ainsi que son contexte de diffusion. Par exemple, la manifestation de la version restaurée de \textit{L'Aigle à deux têtes} décrirait la sortie de cette version en Blu-Ray et DVD chez Rimini Editions. Enfin, le niveau Item est également conservé afin de décrire l’exemplaire physique ou le fichier numérique unique issu de la Manifestation d’une Œuvre ou de sa Variante. Cet exemplaire est concrètement préservé et géré par l’institution et décrit par des métadonnées de conservation.  

Cette restructuration du FRBR permet ainsi de répondre aux besoins métiers spécifiques de gestion de collections d’images animées. Selon la FIAF, la modélisation de cette réadaptation du FRBR au sein d’un outil de gestion des collections devrait permettre aux utilisateurs métier “d’identifier, sélectionner et acquérir” efficacement toutes les variantes d’une œuvre, l’ensemble des manifestations d’une version d’une œuvre, les exemplaires physiques disponibles d’une manifestation d’une œuvre \footcite{zotero-item-646}. La FIAF reprend ainsi les besoins utilisateurs fondateurs de la structuration du FRBR bibliographique en les adaptant à une nécessité de pouvoir circuler entre les différentes versions d’une œuvre à différents stades du processus de l’industrie cinématographique.  

Toutefois, l’implémentation d’un tel modèle au sein d’un outil de gestion des collections se heurte aux spécificités des différents corpus audiovisuels, dont certains ne possèdent pas de manifestation par exemple. En effet, si le modèle s'applique adéquatement aux corpus de cinéma dits "classiques" dont la chaîne de production industrielle est respectée : un film est réalisé à l'issue d'un tournage et d'un montage (Œuvre), puis distribué (Manifestation), et copié en multiples exemplaires (Item). Or, dans le cas de corpus spécifiques, expérimentaux, personnels, ou amateurs, qui font l’objet de descriptions complexes, la modélisation FRBR classique s'avère limitée. Ce type de corpus est notamment peu distribué ou diffusé, rendant notamment le niveau de Manifestation obsolète. 

Ainsi, le FRBR reste un modèle qu’il s’agit de réadapter à chaque contexte institutionnel \footcite{boeuf2008}. Le manuel de la FIAF agit comme un cadre théorique de structuration de la description pour la globalité des archives audiovisuelles, qui subit une redéfinition pour chaque type de corpus audiovisuel : amateur, commercial, documentaire, grand public, etc.  

Reproductibilité technique ? \footcite{zotero-item-669} 